\documentclass{article}

\usepackage[T1]{fontenc}
\usepackage[utf8]{inputenc}
\usepackage[a4paper, margin=2.5cm]{geometry}
\usepackage[protrusion=true,expansion=false]{microtype}
\usepackage{graphicx}
\usepackage{amsmath,amssymb}
\usepackage{booktabs}
\usepackage{setspace}
\usepackage[
  backend=biber,
  style=numeric,
  citestyle=numeric-comp,
  sorting=none,
  maxbibnames=10,
  giveninits=true,
  uniquename=init,
  doi=true,
  url=true,
  isbn=false,
  eprint=false
]{biblatex}
\addbibresource{report.bib}

\usepackage[dvipsnames,svgnames,x11names]{xcolor}
\usepackage[
  colorlinks=true,
  linkcolor=blue!60!black,
  citecolor=blue!60!black,
  urlcolor=blue!60!black,
  pdftitle={Advanced Topics in Intelligent Interactive Systems Report: An Empirical Study of LLM-based Blameworthiness Apportionment in Multi-agent Settings},
  pdfauthor={Rodrigo Cort\'es S\'anchez}
]{hyperref}
\usepackage[capitalise,noabbrev]{cleveref}
\usepackage{caption}
\captionsetup{font=small, labelfont=bf, skip=6pt}

\setlength{\parindent}{1em}
\setlength{\parskip}{0.5em}

\title{\huge An Empirical Study of LLM-based Blameworthiness Apportionment in Multi-agent Settings\\
  \vspace{0.5cm}
  \large Advanced Topics in Intelligent Interactive Systems \\
  \large Independent Studies Report}
\author{Rodrigo Cort\'es S\'anchez}
\date{\today}

\begin{document}

\maketitle

\begin{abstract}
  % Briefly state the motivation, approach, key findings, and significance.
\end{abstract}

\tableofcontents
\newpage

%-----------------------------

\section{Introduction}

%-----------------------------

\subsection{Perspectives on Blame and Responsibility}

% - Open by saying blameworthiness is not one single idea; different traditions emphasize control, alternatives, intention, causality, knowledge, and group structure.
% - Use **Frankfurt** to mark the classic shift away from “could have done otherwise” as the core condition for responsibility.
% - Use **Fischer and Ravizza** to show the control-based view: responsibility depends on actual guidance or ownership of the action.
% - Add **Malle** as the psychological angle: blame is shaped by perceived intent, causality, foreseeability, and preventability.
% - Move to **Chockler and Halpern** for the formal causal account: blame and responsibility can be measured in degrees using structural models.
% - Then **Halpern and Kleiman-Weiner**: blame also depends on the agent’s epistemic state, not just causation.
% - End with **Friedenberg and Halpern**: in multiagent settings, blame must be assigned to groups and then distributed to individuals, since outcomes are jointly produced.
% - The overall narrative is: from philosophical responsibility, to psychological blame, to formal causal and epistemic models, to collective blame allocation.

\subsection{Causal Models and Blameworthiness in Multi-agent Systems}

\subsection{The Committee of Voters Example}

\subsection{Responsibility of AI Systems}

% As AI systems are increasingly embedded into our society, we argue that as a means to support ethical and humancentred AI, it is necessary to develop precise models for reasoning about responsibility of AI systems (Dastani, 2023)

\subsubsection{Large Language Models as Agents}

\subsubsection{Large Language Models as Judges}

\subsubsection{Blameworthiness Apportionment in Multi-agent Settings}

\subsection{Research Questions}

%-----------------------------

\section{Methodology}

\subsection{Replication of the Committee of Voters Example \small (FH Blameworthiness, 2023) }

\subsection{}

\section{Results}

\subsection{}

\section{Discussion}

\subsection{}

\section{Limitations}

\section{Conclusions and Future Work}

% What was replicated and demonstrated
% Future Work
% Closing

\newpage

\printbibliography

\end{document}
